\section{Discussion and Limitations}
\label{sec:discussion}
\textbf{Take-aways.} (i) The value of source-model information depends on whether the model under test resembles the sources \emph{in strength}. Under an IID model split, IRT/regression estimators cut RMSE by 17--29\% at $k=50$ and give better rankings; under extrapolation the same estimators are 15--140\% \emph{worse} than a random-subset mean, with errors that grow with the accuracy gap; with held-out model families but no strength shift they work well again. A validation that splits models at random therefore says little about the intended use (evaluating a newer, better model). (ii) Where the source models are used only to \emph{choose} items, informed deterministic subsets do not reliably beat random sampling with a plain mean (Anchor-IRT: 0.90 IID but 1.06/0.96 under extrapolation; anchors, D-optimal and max-variance are worse than random); only source-difficulty \emph{stratification}, which remains a random sample, consistently does (12--16\% lower RMSE in every regime). (iii) A per-target recalibration (two parameters) removes most of the plug-in's extrapolation bias and brings p-IRT to parity with the random mean (0.96), but no better, and with worse rankings. (iv) A cross-fitted difference estimator makes model assistance safe (unbiased, calibrated intervals), but the variance reduction is modest (10--17\%) since most residual variance is item-level Bernoulli noise; a cohort-adaptive shrinkage step recovers most of the IID gain (RMSE $-25\%$ at $k=50$, $-34\%$ at $k=10$) while staying at the level of CF under extrapolation, at the cost of a small bias. (v) No estimator improved \emph{rankings} under extrapolation by more than 0.005 Spearman at $k\ge50$.

\textbf{Relation to prior findings.} Our observations that random sampling is competitive as a selector and that prediction degrades for models more accurate than the sources agree qualitatively with \citet{pacchiardi2024100} and \citet{zhang2025benchmark}; our contribution is the controlled separation of selection and estimation, the bias diagnostics, the exact-unbiasedness/coverage verification, and the cohort-adaptive shrinkage. We did not run the authors' released code, so we do not claim to reproduce their numbers; our IRT-based baselines are our own implementations following the descriptions in \citet{maiapolo2024tiny} and differ from the released tinyBenchmarks code (e.g.\ our anchor clustering uses an SVD embedding, our IRT fit is a MAP fit with fixed priors).

\textbf{Limitations.}
\emph{Protocol.} Extrapolation is simulated by ranking models by their own full-benchmark accuracy; real shifts (new architectures, training data, families) may differ and no release dates were available. Model families are controlled only in the family hold-out regime (7 benchmarks). The target cohort is evaluated on one shared subset; CF-EB needs such a cohort ($T\gtrsim10$) and is not applicable to a single model.
\emph{Methods.} Baselines are our re-implementations and use a single IRT family; stronger learned predictors or different selectors (e.g.\ adaptive acquisition, \citealt{li2025active}) were not tested, and we did not tune baselines per benchmark. CF is exactly unbiased only for simple random sampling; for the stratified sample the equal-probability formula is approximate (its mean signed error at $k{=}50$ is at most 0.13 points in all regimes, but we did not study its variance estimator or coverage). The recalibration penalties were set in a small pilot on three benchmarks in the Extrap-50 regime, so recalibrated numbers may be mildly optimistic; baselines were otherwise not tuned per benchmark (a weaker ability prior improves p-IRT, \cref{sec:ablate}).
\emph{Evidence.} Benchmarks are few (13) and heterogeneous; the confidence intervals mix between-benchmark and between-split variation, and replicates share the same pool of models, so they do not capture uncertainty about the population of future models. Data are binary correctness from a fixed set of public leaderboards (possibly contaminated or containing models with degenerate outputs); we did not study continuous metrics, multi-benchmark pooling, or cost-based stopping rules.
\emph{Generality.} The relative gains we report (about 10--15\% for safe methods) are modest: unbiased subset estimation cannot go much below the Bernoulli noise of $k$ items, so in extrapolation the practical advice is to use more items or to treat model-based predictions as unreliable unless corrected.
